% GENERATED FILE. Edit canonical lessons, then run npm run generate:books.
% canonical-source-hash: fnv1a64:2194b769f182bb30
% canonical-lessons: MR-C468-mrdu, MR-C468-saumya, MR-C468-halva, MR-C468-samvedansil, MR-C468-paropkari

\chapter{Gentle Natures}
\label{ch:mr-gentle-natures}
\hlchaptermodality{468}

\begin{chapteropening}
\textbf{By the end of this chapter:} \emph{I can describe a gentle, sensitive or helpful person.}

Complete the last lesson of chapter 468: I can describe a gentle, sensitive or helpful person.
\end{chapteropening}

\section[\texorpdfstring{\mr{मृदू} (mṛdū)}{मृदू (mṛdū)}]{\mr{मृदू} (mṛdū) — soft-hearted; gentle}

\label{lesson:MR-C468-mrdu}



\begin{quote}
Before the new one: say the Marathi for to stay alive, then the Marathi for to wobble.
\end{quote}

\subsection*{You'll want to know: \mr{मृदू}}
\textbf{\mr{मृदू}} — \emph{mṛdū} — \textquotedblleft{}soft-hearted; gentle\textquotedblright{}.

\begin{quote}
\mr{आमच्या} \mr{शिक्षिका} \mr{स्वभावाने} \mr{मृदू} \mr{आहेत}.
\end{quote}

\begin{quote}
\emph{āmcyā śikṣikā svabhāvāne mṛdū āhet.}
\end{quote}

\begin{quote}
Our teacher is gentle by nature.
\end{quote}

Unchanging: \mr{मृदू} \mr{आवाज}, a gentle voice. Sanskrit writes the word with a short final vowel.

\begin{cousinweb}
\textbf{\mr{मृदू}} (\emph{mṛdū}) is a tatsama, taken from Sanskrit \emph{mṛdu}, \textquotedblleft{}soft, weak\textquotedblright{} (CDIAL 10292), with the final vowel lengthened as Marathi writes it.
\end{cousinweb}

\begin{grammarlens}[title={What you've built}]
One more word for describing what you lived and felt.
\end{grammarlens}

\subsection*{Your turn}
\begin{itemize}
\raggedright
  \item \emph{Say it:} \emph{mṛdū}
  \item \emph{Say it:} \emph{mṛdū}, once more
  \item \emph{Say it:} \emph{āmcyā śikṣikā svabhāvāne mṛdū āhet.} — Our teacher is gentle by nature.
  \item \emph{Recall:} say the Marathi for to stay alive, then the Marathi for to wobble, then say \emph{mṛdū} again
\end{itemize}

\begin{tcolorbox}[breakable,colback=teal!4,colframe=teal!35!black,title={Before you move on}]
What does \mr{मृदू} mean? (\textquotedblleft{}Soft-hearted; gentle\textquotedblright{}.) Say it once more.
\end{tcolorbox}
\section[\texorpdfstring{\mr{सौम्य} (saumya)}{सौम्य (saumya)}]{\mr{सौम्य} (saumya) — mild}

\label{lesson:MR-C468-saumya}



\begin{quote}
Before the new one: say the Marathi for to wobble, then the Marathi for soft-hearted.
\end{quote}

\subsection*{You'll want to know: \mr{सौम्य}}
\textbf{\mr{सौम्य}} — \emph{saumya} — \textquotedblleft{}mild\textquotedblright{}.

\begin{quote}
\mr{आजोबांचा} \mr{स्वभाव} \mr{सौम्य} \mr{आहे}.
\end{quote}

\begin{quote}
\emph{ājobāñcā svabhāv saumya āhe.}
\end{quote}

\begin{quote}
Grandfather has a mild nature.
\end{quote}

Unchanging: of a nature, a medicine or the weather.

\begin{cousinweb}
\textbf{\mr{सौम्य}} (\emph{saumya}) is a tatsama, taken from Sanskrit \emph{saumya}, \textquotedblleft{}gentle and placid\textquotedblright{} (CDIAL 13620).
\end{cousinweb}

\begin{grammarlens}[title={What you've built}]
One more word for describing what you lived and felt.
\end{grammarlens}

\subsection*{Your turn}
\begin{itemize}
\raggedright
  \item \emph{Say it:} \emph{saumya}
  \item \emph{Say it:} \emph{saumya}, once more
  \item \emph{Say it:} \emph{ājobāñcā svabhāv saumya āhe.} — Grandfather has a mild nature.
  \item \emph{Recall:} say the Marathi for to wobble, then the Marathi for soft-hearted, then say \emph{saumya} again
\end{itemize}

\begin{tcolorbox}[breakable,colback=teal!4,colframe=teal!35!black,title={Before you move on}]
What does \mr{सौम्य} mean? (\textquotedblleft{}Mild\textquotedblright{}.) Say it once more.
\end{tcolorbox}
\section[\texorpdfstring{\mr{हळवा} (haḷvā)}{हळवा (haḷvā)}]{\mr{हळवा} (haḷvā) — tender-hearted}

\label{lesson:MR-C468-halva}



\begin{quote}
Before the new one: say the Marathi for soft-hearted, then the Marathi for mild.
\end{quote}

\subsection*{You'll want to know: \mr{हळवा}}
\textbf{\mr{हळवा}} — \emph{haḷvā} — \textquotedblleft{}tender-hearted\textquotedblright{}.

\begin{quote}
\mr{माझी} \mr{बहीण} \mr{खूप} \mr{हळवी} \mr{आहे}, \mr{लगेच} \mr{रडते}.
\end{quote}

\begin{quote}
\emph{mājhī bahīṇ khūp haḷvī āhe, lagec raḍte.}
\end{quote}

\begin{quote}
My sister is very tender-hearted; she cries at once.
\end{quote}

Agrees: \mr{हळवा}, \mr{हळवी}, \mr{हळवं}. \mr{हळवं} \mr{मन}, a sensitive heart.

\begin{cousinweb}
\textbf{\mr{हळवा}} (\emph{haḷvā}) is usually traced to Sanskrit \emph{laghu}, \textquotedblleft{}light, quick\textquotedblright{}: CDIAL lists Marathi \emph{haḷū}, \textquotedblleft{}light\textquotedblright{}, and \emph{haḷvā}, \textquotedblleft{}ripening quickly\textquotedblright{}, under it (CDIAL 10896). The sense \textquotedblleft{}tender-hearted\textquotedblright{} is not given there, so the link is not certain.
\end{cousinweb}

\begin{grammarlens}[title={What you've built}]
One more word for describing what you lived and felt.
\end{grammarlens}

\subsection*{Your turn}
\begin{itemize}
\raggedright
  \item \emph{Say it:} \emph{haḷvā}
  \item \emph{Say it:} \emph{haḷvā}, once more
  \item \emph{Say it:} \emph{mājhī bahīṇ khūp haḷvī āhe, lagec raḍte.} — My sister is very tender-hearted; she cries at once.
  \item \emph{Recall:} say the Marathi for soft-hearted, then the Marathi for mild, then say \emph{haḷvā} again
\end{itemize}

\begin{tcolorbox}[breakable,colback=teal!4,colframe=teal!35!black,title={Before you move on}]
What does \mr{हळवा} mean? (\textquotedblleft{}Tender-hearted\textquotedblright{}.) Say it once more.
\end{tcolorbox}
\section[\texorpdfstring{\mr{संवेदनशील} (saṁvedanśīl)}{संवेदनशील (saṁvedanśīl)}]{\mr{संवेदनशील} (saṁvedanśīl) — sensitive}

\label{lesson:MR-C468-samvedansil}



\begin{quote}
Before the new one: say the Marathi for mild, then the Marathi for tender-hearted.
\end{quote}

\subsection*{You'll want to know: \mr{संवेदनशील}}
\textbf{\mr{संवेदनशील}} — \emph{saṁvedanśīl} — \textquotedblleft{}sensitive\textquotedblright{}.

\begin{quote}
\mr{संवेदनशील} \mr{माणूस} \mr{लवकर} \mr{दुखावतो}.
\end{quote}

\begin{quote}
\emph{saṁvedanśīl māṇūs lavkar dukhāvto.}
\end{quote}

\begin{quote}
A sensitive person is quickly hurt.
\end{quote}

\mr{संवेदना}, feeling, + \mr{शील}, of the nature of. Unchanging.

\begin{cousinweb}
\textbf{\mr{संवेदनशील}} (\emph{saṁvedanśīl}) is a learned compound built from two Sanskrit words, \emph{saṃvedana}, \textquotedblleft{}perceiving, feeling\textquotedblright{}, and \emph{śīla}, \textquotedblleft{}habit, nature\textquotedblright{} (Monier-Williams).
\end{cousinweb}

\begin{grammarlens}[title={What you've built}]
One more word for describing what you lived and felt.
\end{grammarlens}

\subsection*{Your turn}
\begin{itemize}
\raggedright
  \item \emph{Say it:} \emph{saṁvedanśīl}
  \item \emph{Say it:} \emph{saṁvedanśīl}, once more
  \item \emph{Say it:} \emph{saṁvedanśīl māṇūs lavkar dukhāvto.} — A sensitive person is quickly hurt.
  \item \emph{Recall:} say the Marathi for mild, then the Marathi for tender-hearted, then say \emph{saṁvedanśīl} again
\end{itemize}

\begin{tcolorbox}[breakable,colback=teal!4,colframe=teal!35!black,title={Before you move on}]
What does \mr{संवेदनशील} mean? (\textquotedblleft{}Sensitive\textquotedblright{}.) Say it once more.
\end{tcolorbox}
\section[\texorpdfstring{\mr{परोपकारी} (paropkārī)}{परोपकारी (paropkārī)}]{\mr{परोपकारी} (paropkārī) — helpful to others}

\label{lesson:MR-C468-paropkari}



\begin{quote}
Before the new one: say the Marathi for tender-hearted, then the Marathi for sensitive.
\end{quote}

\subsection*{You'll want to know: \mr{परोपकारी}}
\textbf{\mr{परोपकारी}} — \emph{paropkārī} — \textquotedblleft{}helpful to others\textquotedblright{}.

\begin{quote}
\mr{आमचे} \mr{शेजारी} \mr{परोपकारी} \mr{आहेत}.
\end{quote}

\begin{quote}
\emph{āmce śejārī paropkārī āhet.}
\end{quote}

\begin{quote}
Our neighbours are always ready to help others.
\end{quote}

\mr{पर}, other, + \mr{उपकार}, a kindness: \mr{परोपकार}, doing good to others.

\begin{cousinweb}
\textbf{\mr{परोपकारी}} (\emph{paropkārī}) is a tatsama, taken from Sanskrit \emph{paropakārin}, \textquotedblleft{}assisting others, beneficent\textquotedblright{} (Monier-Williams).
\end{cousinweb}

\begin{grammarlens}[title={What you've built}]
One more word for describing what you lived and felt.
\end{grammarlens}

\subsection*{Your turn}
\begin{itemize}
\raggedright
  \item \emph{Say it:} \emph{paropkārī}
  \item \emph{Say it:} \emph{paropkārī}, once more
  \item \emph{Say it:} \emph{āmce śejārī paropkārī āhet.} — Our neighbours are always ready to help others.
  \item \emph{Recall:} say the Marathi for tender-hearted, then the Marathi for sensitive, then say \emph{paropkārī} again
\end{itemize}

\begin{tcolorbox}[breakable,colback=teal!4,colframe=teal!35!black,title={Before you move on}]
What does \mr{परोपकारी} mean? (\textquotedblleft{}Helpful to others\textquotedblright{}.) Say it once more.
\end{tcolorbox}
